\documentclass[11pt,a4paper]{article}
\usepackage[margin=27mm,headheight=15pt]{geometry}
\usepackage[T1]{fontenc}
\usepackage[utf8]{inputenc}
\usepackage{amsmath,amsthm}
\usepackage{newpxtext,newpxmath}
\usepackage{microtype}
\usepackage{mathtools}
\usepackage{booktabs,array,longtable}
\usepackage{xcolor}
\usepackage{enumitem,aliascnt}
\usepackage{fancyhdr}
\usepackage{titlesec}
\usepackage{listings}
\usepackage[most]{tcolorbox}
\definecolor{OliveInk}{HTML}{46523B}
\definecolor{MutedInk}{HTML}{62675D}
\definecolor{SoftOlive}{HTML}{F1F3EB}
\definecolor{RuleOlive}{HTML}{B6BEA5}
\usepackage[colorlinks=true,linkcolor=OliveInk,citecolor=OliveInk,urlcolor=OliveInk,
 pdfauthor={Research report prepared for Vladimir Reshetnikov},
 pdftitle={A sharp order-five threshold for log-concavity of higher-order Stirling subset numbers}]{hyperref}
\usepackage[nameinlink,noabbrev]{cleveref}
\titleformat{\section}{\normalfont\Large\bfseries\color{OliveInk}}{\thesection}{0.6em}{}
\titleformat{\subsection}{\normalfont\large\bfseries\color{OliveInk}}{\thesubsection}{0.6em}{}
\pagestyle{fancy}
\fancyhf{}
\fancyhead[L]{\small\color{MutedInk}Higher-order Stirling log-concavity}
\fancyhead[R]{\small\color{MutedInk}Research manuscript}
\fancyfoot[C]{\small\thepage}
\renewcommand{\headrulewidth}{0.3pt}
\setlength{\parskip}{0.32em}
\setlength{\parindent}{1.25em}
\setlength{\emergencystretch}{2em}
\setcounter{tocdepth}{1}
\numberwithin{equation}{section}
\newtheorem{theorem}{Theorem}[section]
\newaliascnt{lemma}{theorem}
\newtheorem{lemma}[lemma]{Lemma}
\aliascntresetthe{lemma}
\crefname{lemma}{lemma}{lemmas}
\newaliascnt{proposition}{theorem}
\newtheorem{proposition}[proposition]{Proposition}
\aliascntresetthe{proposition}
\crefname{proposition}{proposition}{propositions}
\newaliascnt{corollary}{theorem}
\newtheorem{corollary}[corollary]{Corollary}
\aliascntresetthe{corollary}
\crefname{corollary}{corollary}{corollaries}
\theoremstyle{definition}
\newaliascnt{definition}{theorem}
\newtheorem{definition}[definition]{Definition}
\aliascntresetthe{definition}
\crefname{definition}{definition}{definitions}
\newaliascnt{example}{theorem}
\newtheorem{example}[example]{Example}
\aliascntresetthe{example}
\crefname{example}{example}{examples}
\theoremstyle{remark}
\newaliascnt{remark}{theorem}
\newtheorem{remark}[remark]{Remark}
\aliascntresetthe{remark}
\crefname{remark}{remark}{remarks}
\newcommand{\Z}{\mathbb Z}
\newcommand{\R}{\mathbb R}
\newcommand{\Q}{\mathbb Q}
\newcommand{\N}{\mathbb N}
\newcommand{\lc}{\operatorname{LC}}
\newcommand{\Snum}[3]{S^{(#1)}_{#2,#3}}
\newcommand{\Cnum}[3]{C^{(#1)}_{#2,#3}}
\newcommand{\assocS}[3]{\genfrac\{\}{0pt}{}{#1}{#2}_{\!\ge #3}}
\newcommand{\assocC}[3]{\genfrac[]{0pt}{}{#1}{#2}_{\!\ge #3}}
\newcommand{\Turan}{Tur\'an}
\lstset{basicstyle=\small\ttfamily,columns=fullflexible,keepspaces=true,
 breaklines=true,frame=single,rulecolor=\color{RuleOlive},
 backgroundcolor=\color{SoftOlive},showstringspaces=false,tabsize=4}
\tcbset{colback=SoftOlive,colframe=RuleOlive,boxrule=0.5pt,arc=1mm,
 left=3mm,right=3mm,top=2mm,bottom=2mm}

\begin{document}
\begin{center}
{\small\scshape\color{MutedInk}Combinatorics \quad / \quad Integer sequences \quad / \quad Recurrences}\par
\vspace{0.6cm}
{\LARGE\bfseries\color{OliveInk}A sharp order-five threshold for\par
log-concavity of higher-order\par Stirling subset numbers}\par
\vspace{0.45cm}
{\large An elementary recurrence proof, explicit inequalities,\par
and applications to Stirling cycle triangles}\par
\vspace{0.45cm}
{\small Research report prepared for Vladimir Reshetnikov\par
19 September 2026}
\end{center}
\vspace{0.3cm}
\begin{abstract}
For an integer $r\ge1$, let $S^{(r)}_{n,k}$ count partitions of an
$(n+(r-1)k)$-element set into $k$ blocks, each of size at least $r$.
We prove that all rows of this diagonal-shifted Stirling triangle are
log-concave if and only if $1\le r\le5$. The positive-support interior
inequalities are strict. This establishes the log-concavity assertion
of Deb--Sokal's Conjecture~1.4(c), including its cases $r=3,4,5$.
The proof reduces a two-term recurrence to a quadratic form. For
$r=5$, a factored discriminant proves positivity and supplies an explicit
lower bound for each \Turan{} difference. The converse follows from an
exact counterexample at $(r,n,k)=(6,4,2)$ and a decreasing factorial
ratio that gives counterexamples in row three for every $r\ge7$.
The same method proves log-concavity for the cycle triangles of orders
three and four and for a continuous family of interpolating recurrences.
We also exhibit a precise obstruction to extending the method to the
fifth-order cycle triangle; that separate conjectural case is not settled
here. All arguments are given in ordinary mathematical form. An accompanying
standard-library Python package checks the polynomial identities exactly
and independently verifies the number arrays by coefficient extraction.
\end{abstract}

\begin{tcolorbox}
\textbf{Scope and verification.} The selected subset-number assertion has a
complete proof below; the claim is not based on a numerical search.
The originating preprint still states it as a conjecture in the version
accessed for this report. A targeted literature check did not locate a later
resolution, but this does not certify priority. This manuscript has not
undergone independent peer review or proof-assistant verification.
\end{tcolorbox}
\newpage
\tableofcontents
\vspace{0.9cm}
\begin{tcolorbox}[title={\bfseries Reading guide},colbacktitle=SoftOlive,
 coltitle=OliveInk]
The proof of the selected conjecture is in Sections~2--6.
Section~3 isolates the local inequality, and Section~4 contains the
order-five factorization. Sections~7--9 give the cycle-number extensions,
an exact obstruction, and an asymptotic explanation. Sections~10--11
supply generating functions, sequence data, and the computational audit.
The appendices collect the polynomial formulas and source-status details.
\end{tcolorbox}
\newpage

\section{The problem and the results}\label{sec:results}

Deb and Sokal introduce higher-order Stirling triangles and conjecture
row log-concavity through order five in
\cite[Conjectures~1.3(c) and 1.4(c)]{DebSokal}.
The cycle-number cases $r=3,4,5$ also appear as Problem~30.9 in the
BCC30 problem collection \cite{Cameron}. We select the subset-number
assertion because its recurrence has a particularly useful linear
coefficient. The second-order instance is the Ward triangle, recorded
as OEIS~A134991 after removing the zero column \cite{OEIS}.

Write $\assocS{N}{k}{r}$ for the number of set partitions of $[N]$ into
$k$ blocks of size at least $r$, and $\assocC{N}{k}{r}$ for the number
of permutations of $[N]$ with $k$ cycles of length at least $r$.
Here $[N]=\{1,\ldots,N\}$ and the empty structure has size zero.
The arrays of interest are
\begin{equation}\label{eq:definitions}
 \Snum r n k=\assocS{n+(r-1)k}{k}{r},\qquad
 \Cnum r n k=\assocC{n+(r-1)k}{k}{r}.
\end{equation}
Their initial row is $(1,0,0,\ldots)$, and entries outside $0\le k\le n$
are zero. In every nonzero row, the positive entries are precisely those
with $1\le k\le n$.

\begin{definition}
A nonnegative sequence $(a_k)$ is \emph{log-concave} if
$a_k^2\ge a_{k-1}a_{k+1}$ at every relevant index. For a finite sequence
we extend it by zeros. We say it is \emph{strictly log-concave on the
interior of its positive support} when the inequality is strict whenever
all three entries are positive.
\end{definition}

\begin{theorem}[Sharp subset-number classification]\label{thm:main}
For a positive integer $r$, the following are equivalent:
\begin{enumerate}[label=(\roman*),nosep]
\item Every row of $S^{(r)}$ is log-concave.
\item $1\le r\le5$.
\end{enumerate}
For these five orders, and for every $n\ge3$ and $2\le k\le n-1$,
\begin{equation}\label{eq:main-strict}
 (\Snum r n k)^2>\Snum r n {k-1}\,\Snum r n {k+1}.
\end{equation}
For $r=6$ the first failing row is $n=4$. For every $r\ge7$ the first
failing row is $n=3$.
\end{theorem}

\begin{theorem}[Cycle-number consequence]\label{thm:cycles}
Every row of $C^{(r)}$ is log-concave for $1\le r\le4$, with strict
inequalities on the interior of the positive support. In particular,
the log-concavity cases $r=3,4$ of Deb--Sokal's Conjecture~1.3(c) hold.
\end{theorem}

These theorems concern the log-concavity clauses only. The same numbered
conjectures in \cite{DebSokal} contain other assertions; no claim about
all their clauses follows from this manuscript. In particular, we do not
settle the order-five cycle case, the general total positivity of the
triangles, or their Hankel-total-positivity questions.

A crucial distinction is that $n$ in \eqref{eq:definitions} is \emph{not}
the cardinality of the underlying set. Moving right in a fixed row changes
that cardinality by $r-1$. Thus a theorem about a fixed-$N$ row of ordinary
associated Stirling numbers would not, by itself, solve this problem.

\subsection{What makes the proof work}

The proof has three ingredients. First, a geometric rescaling removes the
factorial denominator from the subset recurrence. Second, three elementary
consequences of log-concavity reduce the next-row inequality to a quadratic
form in two adjacent entries of the previous row. Finally, the triangular
index restriction bounds its mixed coefficient. Orders three and four then
follow from nonnegative coefficients. At order five the mixed coefficient
can be negative, so it must be controlled rather than discarded.

This is closely related to Sagan's recurrence criterion
\cite[Theorem~1]{Sagan}, whose nonnegative mixed-term hypothesis is sufficient
for some of the lower-order cases. Our local calculation retains the
positive square terms and allows a negative mixed term. We do not claim
that the general quadratic-form observation is historically new; the
specific application and factorization are what establish the selected
assertion here. A more recent paper of Shankar \cite{Shankar} studies
other sufficient conditions for triangular recurrences. Its principal
power-of-affine-coefficients family is different from the falling-factorial
weight used below.

\section{Recurrences, normalization, and boundary values}\label{sec:recurrence}

Set $d=r-1$, and define the falling-factorial polynomial
\begin{equation}\label{eq:P}
 P_d(t)=\prod_{j=0}^{d-1}(t-j),\qquad P_0(t)=1.
\end{equation}
The two recurrences in \cite[Lemma~1.1]{DebSokal} can be derived directly
from the distinguished largest label, as follows.

\begin{proposition}\label{prop:recurrences}
For $n\ge1$, $1\le k\le n$, and $t=n+dk-1$,
\begin{align}
 \Snum r n k
   &=\binom{t}{d}\Snum r {n-1}{k-1}+k\Snum r {n-1}k,
       \label{eq:Srec}\\
 \Cnum r n k
   &=P_d(t)\Cnum r {n-1}{k-1}+t\Cnum r {n-1}k.
       \label{eq:Crec}
\end{align}
\end{proposition}
\begin{proof}
Let $N=n+dk$. In a partition counted by $\assocS Nk r$, the block
containing $N$ either has exactly $r$ elements or has more than $r$.
In the first case choose its $r-1$ other elements and partition the remaining
$N-r$ elements into $k-1$ admissible blocks. In the second case remove $N$;
there is an admissible partition of $N-1$ elements and a choice of one of its
$k$ blocks in which to reinsert $N$. Hence
\[
 \assocS Nk r=\binom{N-1}{r-1}\assocS{N-r}{k-1}r
                 +k\assocS{N-1}kr.
\]
For permutations, a cycle of size exactly $r$ has $(r-1)!$ cyclic
orderings once its elements have been chosen. Otherwise, removing $N$
leaves a permutation into whose cycles $N$ can be reinserted after any
of the $N-1$ remaining labels. The corresponding weights are therefore
$(r-1)!\binom{N-1}{r-1}$ and $N-1$.
Finally,
$N-r=(n-1)+d(k-1)$ and $N-1=(n-1)+dk$, giving the stated recurrences.
\end{proof}

For the subset array define
\begin{equation}\label{eq:normalization}
 T^{(r)}_{n,k}=(d!)^k\Snum r n k.
\end{equation}
Then \eqref{eq:Srec} becomes
\begin{equation}\label{eq:Trec}
 T^{(r)}_{n,k}=P_d(n+dk-1)T^{(r)}_{n-1,k-1}
                      +kT^{(r)}_{n-1,k}.
\end{equation}
This normalization leaves log-concavity unchanged, since
\begin{equation}\label{eq:scale-difference}
 (T^{(r)}_{n,k})^2-T^{(r)}_{n,k-1}T^{(r)}_{n,k+1}
  =(d!)^{2k}\bigl[(\Snum r n k)^2-\Snum r n{k-1}\Snum r n{k+1}\bigr].
\end{equation}
The cycle array already has the same first weight $P_d(t)$; only its
second weight differs.

For later checks, direct counting gives, for $n\ge1$,
\begin{align}
 \Snum r n1&=1,&
 \Snum r nn&=\frac{(rn)!}{n!(r!)^n},\label{eq:Sboundary}\\
 \Cnum r n1&=(n+r-2)!,&
 \Cnum r nn&=\frac{(rn)!}{n!r^n}.\label{eq:Cboundary}
\end{align}
In particular $C^{(r)}_{1,1}=(r-1)!$; for $r>2$ this is not one.
These values also make the support assertions immediate from the recurrences.
The rows $n=0,1,2$ have at most two positive entries, so their
log-concavity is automatic.

\section{A sharp local quadratic-form criterion}\label{sec:criterion}

Let $(x_j)$ be a nonnegative log-concave sequence with no internal zeros.
At a fixed index $k$ suppose $x_{k-1}>0$ and $x_k>0$, and put
\[
 w=x_{k-2},\qquad u=x_{k-1},\qquad v=x_k,\qquad z=x_{k+1}.
\]
Log-concavity implies
\begin{equation}\label{eq:three-slacks}
 wv\le u^2,\qquad uz\le v^2,\qquad wz\le uv.
\end{equation}
The third inequality follows by multiplying the first two and dividing by
$uv>0$. It remains true if $w$ or $z$ is zero.

Consider nonnegative coefficients and the local transform
\begin{equation}\label{eq:local-transform}
 y_j=a_jx_{j-1}+b_jx_j\qquad (j=k-1,k,k+1).
\end{equation}
Use the abbreviations $a=a_k$, $a_-=a_{k-1}$, $a_+=a_{k+1}$ and similarly
for $b$. Define
\begin{equation}\label{eq:ABH}
 A=a^2-a_-a_+,\qquad B=b^2-b_-b_+,\qquad
 H=2ab-a_-b_+-b_-a_+.
\end{equation}

\begin{lemma}[Exact slack decomposition]\label{lem:slacks}
With the preceding notation,
\begin{align}
 y_k^2-y_{k-1}y_{k+1}
  ={}& Au^2+Huv+Bv^2 \notag\\
    &+a_-a_+(u^2-wv)+b_-b_+(v^2-uz)
       +a_-b_+(uv-wz).\label{eq:slack-identity}
\end{align}
Consequently,
\begin{equation}\label{eq:quadratic-bound}
 y_k^2-y_{k-1}y_{k+1}\ge Au^2+Huv+Bv^2.
\end{equation}
\end{lemma}
\begin{proof}
Expand $(au+bv)^2-(a_-w+b_-u)(a_+v+b_+z)$ and collect terms.
This gives \eqref{eq:slack-identity} exactly. Each of the three remaining
brackets is nonnegative by \eqref{eq:three-slacks}, and each multiplier
is nonnegative.
\end{proof}

\begin{proposition}[Local preservation criterion]\label{prop:copositive}
The quadratic form $Au^2+Huv+Bv^2$ is nonnegative for every $u,v>0$
if and only if
\begin{equation}\label{eq:copositive}
 A\ge0,\qquad B\ge0,\qquad H\ge-2\sqrt{AB}.
\end{equation}
Thus these conditions imply the next-row log-concavity inequality.
They are also necessary if the transform is required to satisfy that
inequality for every positive log-concave four-term input.
\end{proposition}
\begin{proof}
If \eqref{eq:copositive} holds, write
\[
 Au^2+Huv+Bv^2=(\sqrt A\,u-\sqrt B\,v)^2
                         +(H+2\sqrt{AB})uv\ge0.
\]
Conversely, sending $v/u$ to zero or infinity forces $A\ge0$ and $B\ge0$.
For $A,B>0$, choosing $v/u=\sqrt{A/B}$ forces the third condition.
If $A=0$ or $B=0$, the same condition, now $H\ge0$, follows by taking
an appropriately small or large positive ratio.

For necessity as a preservation criterion, take the four input values
\[
 (w,u,v,z)=(u^2/v,\ u,\ v,\ v^2/u).
\]
They form a positive geometric progression, so all three slacks in
\eqref{eq:slack-identity} vanish. The next-row difference is exactly the
quadratic form. Extending this four-term sequence by zeros produces
no difficulty at the support boundaries.
\end{proof}

Only nonnegativity on the positive quadrant is needed, not positive
semidefiniteness on all of $\R^2$. In matrix terminology this is
copositivity of a $2\times2$ symmetric matrix. For the hard subset case we
will use the stronger, more easily certified condition
$4AB-H^2>0$ after replacing $H$ by a suitable lower bound.

\begin{remark}[Why a negative mixed term is not fatal]
Sagan's sufficient condition corresponds to $A,B,H\ge0$.
If $H<0$, rejecting the recurrence at once loses the positive contributions
$Au^2$ and $Bv^2$. The exact square-completion threshold in
\eqref{eq:copositive} is what allows us to handle the fifth-order subset
array and the fourth-order cycle array.
\end{remark}

\section{Proof of subset log-concavity through order five}\label{sec:subset-proof}

Fix an interior index $2\le k\le n-1$ in row $n\ge3$, and set
\begin{equation}\label{eq:domain}
 d=r-1,\qquad t=n+dk-1.
\end{equation}
For \eqref{eq:Trec}, the local weights are
\[
 a=P_d(t),\quad a_-=P_d(t-d),\quad a_+=P_d(t+d),\qquad
 b=k,\quad b_-=k-1,\quad b_+=k+1.
\]
Notice that the step in $t$ between \emph{adjacent columns of one row} is
$d$, not $d+1$. The latter number enters a different way: since $n\ge k+1$,
\begin{equation}\label{eq:triangular-bound}
 t\ge(d+1)k,\qquad k\le\frac{t}{d+1},\qquad t\ge2(d+1).
\end{equation}

Let
\begin{align}
 A_d(t)&=P_d(t)^2-P_d(t-d)P_d(t+d),\label{eq:Ad}\\
 H_d(t,k)&=2kP_d(t)-(k+1)P_d(t-d)-(k-1)P_d(t+d).
                                                       \label{eq:Hd}
\end{align}
The remaining quadratic coefficient is simply $B=1$.

\subsection{Two elementary properties of the falling factorial}

\begin{lemma}\label{lem:falling}
For $d\ge1$ and $t>2d-1$, $A_d(t)>0$. For $d\ge2$ in the same domain,
\[
 \Delta_d(t):=P_d(t-d)+P_d(t+d)-2P_d(t)>0.
\]
For $d=0,1$, the centered second difference is zero.
\end{lemma}
\begin{proof}
For $d\ge1$,
\[
 P_d(t-d)P_d(t+d)=\prod_{j=0}^{d-1}\bigl((t-j)^2-d^2\bigr).
\]
Every factor is positive and strictly smaller than $(t-j)^2$, because
$t>2d-1$. Comparing products proves $A_d(t)>0$.

For $d\ge2$, $P_d''(s)>0$ whenever $s>d-1$: differentiate the product
twice, obtaining a sum of positive products of the remaining $d-2$
factors. Thus $P_d$ is strictly convex on the interval $[t-d,t+d]$,
and its centered second difference is positive. The polynomials
$P_0=1$ and $P_1=t$ have zero centered second difference.
\end{proof}

Rewrite the mixed coefficient as
\begin{equation}\label{eq:Hdecreasing}
 H_d(t,k)=P_d(t+d)-P_d(t-d)-k\Delta_d(t).
\end{equation}
For fixed $t$ it is nonincreasing in $k$. Hence \eqref{eq:triangular-bound}
yields the crucial lower bound
\begin{equation}\label{eq:Lbound}
 H_d(t,k)\ge L_d(t):=H_d\!\left(t,\frac{t}{d+1}\right).
\end{equation}
The argument $t/(d+1)$ need not be an integer; $H_d(t,k)$ is affine in $k$,
so this substitution is legitimate.

\subsection{The orders $r=1,2,3,4$}

Direct expansion gives
\begin{center}
\renewcommand{\arraystretch}{1.45}
\begin{tabular}{ccl}
\toprule
$r$ & $d$ & $L_d(t)$\\
\midrule
1 & 0 & $0$\\
2 & 1 & $2$\\
3 & 2 & $\dfrac43(4t-3)$\\
4 & 3 & $\dfrac32(3t^2-15t+44)$\\
\bottomrule
\end{tabular}
\end{center}
For $d=0$, $A_0=0$, $H_0=0$, and the quadratic lower bound is $v^2>0$.
For $d=1$, $A_1=1$ and $H_1=2$, so it is $(u+v)^2>0$.
For $d=2$, the mixed coefficient is positive because $t\ge6$.
For $d=3$, positivity follows even without the domain restriction from
\[
 3t^2-15t+44=3(t-5/2)^2+101/4>0.
\]
In the latter two cases $A_d>0$ by \cref{lem:falling}, so the quadratic
form is strictly positive for $u,v>0$.

\subsection{The order $r=5$}

Here $d=4$, $t\ge10$, and
\begin{align}
 P_4(t)&=t(t-1)(t-2)(t-3),\label{eq:P4}\\
 A_4(t)&=32\bigl(2t^6-18t^5+17t^4+168t^3
                      -55t^2-834t-630\bigr),\label{eq:A4}\\
 L_4(t)&=-\frac{16}{5}\bigl(2t^3+9t^2-161t+255\bigr).
                                                       \label{eq:L4}
\end{align}
The decisive identity is
\begin{equation}\label{eq:main-factor}
\boxed{\begin{aligned}
 4A_4(t)-L_4(t)^2
   =\frac{384}{25}(t-5)(2t^2-6t+9)\,
                    (7t^3-31t^2-126t+1080).
\end{aligned}}
\end{equation}
It is a polynomial identity, not an approximation. Its right-hand side is
strictly positive for every $t\ge10$. Indeed, $t-5>0$,
\[
 2t^2-6t+9=2(t-3/2)^2+9/2>0,
\]
and, on writing $t=10+s$ with $s\ge0$, the final factor is
\begin{equation}\label{eq:positive-cubic}
 7s^3+179s^2+1354s+3720>0.
\end{equation}
For complete auditability, the factorization and all expansions are
checked by coefficient equality in the accompanying \texttt{certificates.py}.
They can also be verified by ordinary multiplication.

By \eqref{eq:Lbound}, for $u,v>0$,
\begin{align}
 A_4(t)u^2+H_4(t,k)uv+v^2
 &\ge A_4(t)u^2+L_4(t)uv+v^2\notag\\
 &=\left(v+\frac{L_4(t)}2u\right)^2
       +\frac{4A_4(t)-L_4(t)^2}{4}\,u^2
 >0.\label{eq:hard-square}
\end{align}
Thus the negative mixed term is more than compensated for by the two
square terms.

\subsection{Completing the induction, including the edges}

The rows $n=0,1,2$ are log-concave. Assume row $n-1$ is log-concave.
For $2\le k\le n-1$, both $u=T^{(r)}_{n-1,k-1}$ and
$v=T^{(r)}_{n-1,k}$ are positive. At $k=2$, $w=0$; at $k=n-1$, $z=0$.
Neither zero causes a problem in \cref{lem:slacks}. The preceding calculations
show that its quadratic lower bound is strictly positive for every order
$r\le5$. Hence all the positive-support interior inequalities in row $n$
are strict. At either endpoint of the positive support one neighboring
entry is zero, and outside the support the inequalities are trivial.
This completes the induction for $T^{(r)}$. Equation
\eqref{eq:scale-difference} transfers the result to $S^{(r)}$.

We have proved the sufficiency direction of \cref{thm:main}, including the strict
inequalities, without using real-rootedness or a finite verification bound.

\section{An explicit fifth-order \Turan{} inequality}\label{sec:quantitative}

The square-completion proof gives more than the sign of a difference.
Let
\begin{equation}\label{eq:Theta}
 \Theta(t)=(t-5)(2t^2-6t+9)(7t^3-31t^2-126t+1080).
\end{equation}

\begin{corollary}\label{cor:quantitative}
For $n\ge3$, $2\le k\le n-1$, and $t=n+4k-1$,
\begin{equation}\label{eq:quantitative}
 (\Snum5nk)^2-\Snum5n{k-1}\Snum5n{k+1}
 \ge\frac{\Theta(t)}{150}\,(\Snum5{n-1}{k-1})^2>0.
\end{equation}
\end{corollary}
\begin{proof}
Drop the nonnegative square in \eqref{eq:hard-square} and the nonnegative
slacks in \eqref{eq:slack-identity}. By \eqref{eq:main-factor}, the difference
for the normalized row is at least
$(96/25)\Theta(t)(T^{(5)}_{n-1,k-1})^2$.
Now divide by $24^{2k}$ and use
$T^{(5)}_{n-1,k-1}=24^{k-1}S^{(5)}_{n-1,k-1}$.
The remaining constant is $96/(25\cdot24^2)=1/150$.
\end{proof}

For instance, the nonzero portion of row three is
\[
 (S^{(5)}_{3,1},S^{(5)}_{3,2},S^{(5)}_{3,3})=(1,462,126126).
\]
Its middle difference is $462^2-126126=87318$.
The lower bound in \eqref{eq:quantitative}, with $t=10$, is $18476$.
It is not asserted to be optimal; its point is to make strict positivity
uniform and explicit.

The result also has immediate shape consequences. In a positive row,
\eqref{eq:main-strict} is equivalent to strict decrease of successive
ratios:
\begin{equation}\label{eq:ratio-monotonicity}
 \frac{S^{(r)}_{n,k}}{S^{(r)}_{n,k-1}}
 >\frac{S^{(r)}_{n,k+1}}{S^{(r)}_{n,k}}.
\end{equation}
Therefore every such row is unimodal. There is either a unique maximum or
two adjacent equal maxima; strict log-concavity does not exclude the latter.
Multiplying entries by $q^k$, for any $q>0$, preserves the theorem.
Normalizing these weighted entries by their sum gives a log-concave
probability mass function on the column index. This is a distribution over
a row of the shifted array, not over partitions of a single fixed set size.

\section{Why order five is the exact cutoff}\label{sec:sharpness}

\subsection{The sixth-order counterexample}

In row four of $S^{(6)}$,
\begin{align*}
 S^{(6)}_{4,1}&=1,\\
 S^{(6)}_{4,2}&=\binom{14}{6}+\frac12\binom{14}{7}=4719,\\
 S^{(6)}_{4,3}&=\frac{19!}{2!(6!)^2\,7!}=23279256.
\end{align*}
The second entry counts block-size patterns $(6,8)$ and $(7,7)$; the third
counts $(6,6,7)$. Consequently,
\begin{equation}\label{eq:r6-counterexample}
 (S^{(6)}_{4,2})^2-S^{(6)}_{4,1}S^{(6)}_{4,3}
 =4719^2-23279256=-1010295<0.
\end{equation}
Row three has positive part $(1,1716,2858856)$ and
$1716^2-2858856=85800>0$. Thus row four is the first failure for $r=6$.

\subsection{A row-three counterexample for every $r\ge7$}

The nonzero part of row three has a uniform closed form:
\begin{equation}\label{eq:row3}
 S^{(r)}_{3,1}=1,\qquad
 S^{(r)}_{3,2}=\binom{2r+1}{r},\qquad
 S^{(r)}_{3,3}=\frac{(3r)!}{6(r!)^3}.
\end{equation}
The two-block case has block sizes $r,r+1$, and the three-block case
has three blocks of size $r$.
Define the ratio that decides log-concavity in this row:
\begin{equation}\label{eq:Fr}
 F(r)=\frac{\binom{2r+1}{r}^{\!2}}{(3r)!/[6(r!)^3]}.
\end{equation}
Cancellation of consecutive factorials gives
\begin{equation}\label{eq:F-ratio}
 \frac{F(r+1)}{F(r)}
  =\frac{4(r+1)^2(2r+3)^2}{3(r+2)^2(3r+2)(3r+1)}.
\end{equation}
The denominator minus the numerator is
\begin{equation}\label{eq:ratio-gap}
 11r^4+55r^3+74r^2+12r-12>0\qquad(r\ge1).
\end{equation}
For a coefficientwise positivity certificate, substitute $r=1+s$; the
polynomial becomes
\[
 11s^4+99s^3+305s^2+369s+140.
\]
Thus $F(r)$ is strictly decreasing on the positive integers. At $r=7$,
\[
 \binom{15}{7}^2=6435^2=41409225<66512160
             =\frac{21!}{6(7!)^3}.
\]
Therefore $F(r)<1$ for every $r\ge7$, proving a failure in row three.
No earlier row can fail because it has at most two positive entries.
Together with \eqref{eq:r6-counterexample}, this proves necessity and
completes \cref{thm:main}.

\section{Cycle triangles and an interpolating family}\label{sec:cycle-proof}

For the cycle recurrence, the first quadratic coefficient is the same
$A_d(t)$, but
\[
 b=t,\quad b_-=t-d,\quad b_+=t+d,\quad B=d^2.
\]
The mixed coefficient is
\begin{equation}\label{eq:Md}
 M_d(t)=2tP_d(t)-(t+d)P_d(t-d)-(t-d)P_d(t+d).
\end{equation}
For $d=1,2,3$ it is
\begin{align}
 M_1(t)&=2,\label{eq:M1}\\
 M_2(t)&=8(t-1),\label{eq:M2}\\
 M_3(t)&=-18(3t-11).\label{eq:M3}
\end{align}
The first two cases have a positive mixed coefficient. For $d=3$, use
\begin{equation}\label{eq:cycle4-discriminant}
 4\cdot9A_3(t)-M_3(t)^2
      =972(t-3)^2(t-1)(t+3)>0\qquad(t\ge8).
\end{equation}
Since $A_3(t)>0$ and $B=9$, the quadratic form is positive definite.
This proves strict log-concavity for orders $r=2,3,4$ by the same induction
and boundary argument as in \cref{sec:subset-proof}.

For completeness, ordinary cycle numbers ($r=1$, $d=0$) can also be
handled without a real-rootedness theorem. The recurrence has $a=1$ and
$b=n-1$, independent of $k$, so $A=B=M=0$. The first slack in
\eqref{eq:slack-identity} nevertheless supplies strictness. For $k=2$,
$w=0$ and $u>0$, hence $u^2-wv>0$. For $k>2$, strict log-concavity in
row $n-1$ at index $k-1$ gives the same inequality. The base row with two
positive entries is automatic. This completes \cref{thm:cycles}.

\subsection{A continuous extension}

The subset and cycle results of orders at most four fit into a single
family. The coefficients below need not be integers.

\begin{theorem}\label{thm:interpolation}
Let $1\le d\le3$, let $\rho,\sigma\ge0$ with $\rho+\sigma>0$, and define
an array $U$ with initial row $U_{0,k}=\delta_{0k}$, zero entries outside
$0\le k\le n$, and recurrence
\begin{equation}\label{eq:interpolating}
 U_{n,k}=P_d(n+dk-1)U_{n-1,k-1}
       +\bigl[\rho(n+dk-1)+\sigma k\bigr]U_{n-1,k}
\end{equation}
for $1\le k\le n$, with $U_{n,0}=0$ for $n\ge1$.
Every row is strictly log-concave on the interior of its positive support.
\end{theorem}
\begin{proof}
All entries on $1\le k\le n$ are positive. At a fixed interior index,
write $t=n+dk-1$. The first quadratic coefficient remains $A_d(t)$.
Since the second weight is affine in $k$ with slope $\rho d+\sigma$,
\[
 B=(\rho d+\sigma)^2,\qquad
 H=\rho M_d(t)+\sigma H_d(t,k).
\]
Our preceding calculations give
\[
 M_d(t)>-2d\sqrt{A_d(t)},\qquad
 H_d(t,k)>-2\sqrt{A_d(t)}
\]
for $d=1,2,3$ in the triangular interior. Multiplying by $\rho,\sigma$
and adding yields
\[
 H>-2(\rho d+\sigma)\sqrt{A_d(t)}=-2\sqrt{A_d(t)B}.
\]
The local criterion proves strict positivity of the quadratic form, and
induction completes the proof.
\end{proof}

The choice $(\rho,\sigma)=(0,1)$ is the normalized subset array, and
$(1,0)$ is the cycle array. This theorem is a genuine parameter extension,
not a convex-combination assertion about already computed rows: the
parameters are inserted in the recurrence at every step.

\section{A rigorous obstruction in the fifth-order cycle case}\label{sec:obstruction}

The preceding method does not automatically settle the remaining cycle
case. In fact, universal preservation of log-concavity is false for its
row operator. At $d=4$,
\begin{equation}\label{eq:M4}
 M_4(t)=-32(t-2)(2t^2+4t-51),
\end{equation}
and
\begin{equation}\label{eq:cycle5-discriminant}
 \boxed{64A_4(t)-M_4(t)^2
       =-9216(t-4)(t+4)(2t-9)(2t^2-6t+9).}
\end{equation}
For $t\ge10$ the right-hand side is negative and $M_4(t)<0$.
Thus the necessary local preservation condition fails.

Here is an exact witness that also respects the support of a preceding
triangular row. Let $n=5$, $q=7350$, and take the artificial input
\[
 (x_0,x_1,x_2,x_3,x_4,x_5)=(0,1,q,q^2,q^3,0).
\]
It is log-concave. Apply the fifth-order cycle row operator
\[
 y_j=P_4(4+4j)x_{j-1}+(4+4j)x_j\qquad(1\le j\le5).
\]
The three relevant output entries are
\begin{equation}\label{eq:witness-values}
 y_2=100080,\qquad y_3=1185408000,\qquad y_4=14223043800000,
\end{equation}
and their \Turan{} difference is
\begin{equation}\label{eq:witness-negative}
 y_3^2-y_2y_4=-18250097040000000<0.
\end{equation}
The choice of $q$ is not arbitrary. At this index $t=16$, the quadratic
form $A_4u^2+M_4uv+16v^2$ is minimized as a function of $v/u$ at
\[
 \frac vu=-\frac{M_4(16)}{32}
       =(16-2)\bigl(2\cdot16^2+4\cdot16-51\bigr)=7350.
\]
The geometric input makes every slack zero and realizes the negative
quadratic value exactly.

\begin{tcolorbox}
\textbf{What this does and does not show.}
The artificial input is not a row of $C^{(5)}$.
Equation~\eqref{eq:witness-negative} disproves only the claim that this
operator preserves \emph{all} log-concave inputs. It does not disprove
log-concavity of the actual fifth-order Stirling cycle rows.
That case remains unresolved by this report.
\end{tcolorbox}

There is a concrete next invariant to investigate. For the actual previous
row, retain the three nonnegative slacks in \eqref{eq:slack-identity}
and seek a lower bound on their weighted sum that dominates the possible
negative part of $A_4u^2+M_4uv+16v^2$. Alternatively, prove a restriction
on the actual adjacent ratio $v/u$ that excludes the dangerous interval
between the two positive roots of
\[
 16q^2+M_4(t)q+A_4(t)=0.
\]
Either approach would use information specific to the array, which
ordinary row log-concavity alone cannot supply. These are proposed
extensions, not additional proved statements.

\section{A leading-order explanation of the cutoff}\label{sec:asymptotics}

The exact proof is polynomial, but its structure has a useful asymptotic
interpretation. Fix $d\ge1$ and let $t\to\infty$. Since
$P_d(t)=t^d+O(t^{d-1})$, centered Taylor expansions give
\begin{align}
 P_d(t+d)-P_d(t-d)&=2d^2t^{d-1}+O(t^{d-2}),\label{eq:asym-diff}\\
 \Delta_d(t)&=d^3(d-1)t^{d-2}+O(t^{d-3}).\label{eq:asym-second}
\end{align}
For $d=1$ the centered second difference is exactly zero, as this formula
permits. Comparing the leading terms in $P_d^2-P_d(t-d)P_d(t+d)$ gives
\begin{equation}\label{eq:asym-A}
 A_d(t)=d^3t^{2d-2}+O(t^{2d-3}).
\end{equation}
For example, this also follows from
$\log P_d(t-d)+\log P_d(t+d)-2\log P_d(t)
=-d^3/t^2+O(t^{-3})$.
Combining \eqref{eq:asym-diff}--\eqref{eq:asym-second} with
$L_d=P_d(t+d)-P_d(t-d)-t\Delta_d/(d+1)$ yields
\begin{equation}\label{eq:asym-L}
 L_d(t)=\frac{d^2(-d^2+3d+2)}{d+1}\,t^{d-1}+O(t^{d-2}).
\end{equation}
Hence the normalized mixed coefficient has limit
\begin{equation}\label{eq:asym-subset-ratio}
 \frac{L_d(t)}{2\sqrt{A_d(t)}}
 \longrightarrow \frac{\sqrt d\,(-d^2+3d+2)}{2(d+1)}.
\end{equation}
At $d=4$ this is $-2/5$, safely above the local threshold $-1$.
At $d=5$ it is $-2\sqrt5/3<-1$. Thus the elementary preservation
mechanism naturally changes between orders five and six. The exact
counterexamples in \cref{sec:sharpness}, rather than this asymptotic
observation alone, prove the cutoff for the actual arrays.

For cycles, \eqref{eq:Md} can be written as
$M_d=d(P_d(t+d)-P_d(t-d))-t\Delta_d$. Therefore
\begin{equation}\label{eq:asym-M}
 M_d(t)=d^3(3-d)t^{d-1}+O(t^{d-2}),\qquad
 \frac{M_d(t)}{2d\sqrt{A_d(t)}}\longrightarrow
                         \frac{\sqrt d(3-d)}2.
\end{equation}
At $d=4$ the latter limit is exactly $-1$, on the boundary.
The leading terms of $64A_4$ and $M_4^2$ cancel; their difference begins
with $-36864t^5$. This explains why the fifth-order cycle operator is
more delicate than its subset counterpart, and why keeping lower-order
terms is essential.

\section{Generating functions and reproducible sequence data}\label{sec:gf}

\subsection{An independent coefficient formula}

Define formal power series
\begin{equation}\label{eq:fg}
 f_r(z)=\sum_{j\ge0}\frac{z^j}{(r+j)!},\qquad
 g_r(z)=\sum_{j\ge0}\frac{z^j}{r+j}.
\end{equation}
For each fixed $k$, counting ordered collections of $k$ blocks and then
dividing by $k!$ gives
\begin{align}
 \sum_{N\ge0}\assocS Nkr\frac{z^N}{N!}
  &=\frac1{k!}\left(\sum_{m\ge r}\frac{z^m}{m!}\right)^k,
                                                       \label{eq:associated-S-egf}\\
 \sum_{N\ge0}\assocC Nkr\frac{z^N}{N!}
  &=\frac1{k!}\left(\sum_{m\ge r}\frac{z^m}{m}\right)^k.
                                                       \label{eq:associated-C-egf}
\end{align}
For cycles, the coefficient $1/m$ is $(m-1)!/m!$.
Substituting $N=n+(r-1)k$ and factoring out $z^{rk}$ gives
\begin{align}
 \Snum rnk&=\frac{(n+(r-1)k)!}{k!}[z^{n-k}]f_r(z)^k,
                                                       \label{eq:S-coefficient}\\
 \Cnum rnk&=\frac{(n+(r-1)k)!}{k!}[z^{n-k}]g_r(z)^k.
                                                       \label{eq:C-coefficient}
\end{align}
These formulas are useful independent checks because they do not use the
two-term recurrence. Only degrees through $n-k$ need be retained.
They reproduce the first-column and diagonal formulas immediately.

\subsection{Row-polynomial operators}

Let $\vartheta=x\,d/dx$, and write
$T^{(r)}_n(x)=\sum_kT^{(r)}_{n,k}x^k$ and
$C^{(r)}_n(x)=\sum_kC^{(r)}_{n,k}x^k$.
Summing the recurrences, with multiplication by $x$ performed \emph{after}
the differential operator, gives
\begin{align}
 T^{(r)}_n(x)
 &=xP_d(n+d-1+d\vartheta)T^{(r)}_{n-1}(x)
                         +\vartheta T^{(r)}_{n-1}(x),\label{eq:Toperator}\\
 C^{(r)}_n(x)
 &=xP_d(n+d-1+d\vartheta)C^{(r)}_{n-1}(x)
                    +(n-1+d\vartheta)C^{(r)}_{n-1}(x).
                                                       \label{eq:Coperator}
\end{align}
Here a polynomial in $\vartheta$ acts diagonally on monomials:
$P_d(n+d-1+d\vartheta)x^j=P_d(n+d-1+dj)x^j$.
Thus the order and indexing of the operator are unambiguous.
These formulas offer a compact generating-polynomial description, but the
proof above does not require an analysis of polynomial zeros.

\subsection{Small rows}

Table~\ref{tab:rows} displays a few of the rows covered by the conjectural
cases. Empty places denote entries outside the triangle, not omitted data.
The zero column has been suppressed.

\begin{table}[htbp]
\centering
\small
\renewcommand{\arraystretch}{1.2}
\begin{tabular}{ccrrrr}
\toprule
$r$ & $n$ & $S_{n,1}^{(r)}$ & $S_{n,2}^{(r)}$ & $S_{n,3}^{(r)}$ & $S_{n,4}^{(r)}$\\
\midrule
3 & 2 & 1 & 10 & &\\
3 & 3 & 1 & 35 & 280 &\\
3 & 4 & 1 & 91 & 2100 & 15400\\
\addlinespace
4 & 2 & 1 & 35 & &\\
4 & 3 & 1 & 126 & 5775 &\\
4 & 4 & 1 & 336 & 45045 & 2627625\\
\addlinespace
5 & 2 & 1 & 126 & &\\
5 & 3 & 1 & 462 & 126126 &\\
5 & 4 & 1 & 1254 & 1009008 & 488864376\\
\bottomrule
\end{tabular}
\caption{Higher-order subset rows, computed exactly from \eqref{eq:Srec}
and checked independently using \eqref{eq:S-coefficient}.}\label{tab:rows}
\end{table}

For $r=3$, the first six subset row sums, starting at $n=0$, are
\[
 1,\ 1,\ 11,\ 316,\ 17592,\ 1612206.
\]
For $r=4$ they are
\[
 1,\ 1,\ 36,\ 5902,\ 2673007,\ 2582136349,
\]
and for $r=5$,
\[
 1,\ 1,\ 127,\ 126589,\ 489874639,\ 5201522272948.
\]
The archive includes row sums through $n=30$ and complete small triangles
through $n=12$ for both kinds and $1\le r\le8$. These are reproducible
integer-sequence data; no assertion is made about whether each sequence
already has a separate OEIS entry.

\section{Exact computational audit}\label{sec:verification}

The mathematical proof is finite algebra plus induction over all rows.
The computations serve two distinct purposes: checking those finite
algebraic certificates, and testing for implementation or indexing mistakes
in the number arrays. Neither numerical root finding nor floating-point
inequalities are used.

\subsection{Polynomial certificates}

The script \texttt{certificates.py} implements elementary polynomial
addition, multiplication, scalar division, and substitution over $\Q$.
Polynomials are tuples of rational coefficients in ascending degree.
All identities are checked by exact equality of the coefficient tuples,
not by evaluation at finitely many points. The recorded run verified
45 identities or shifted-coefficient positivity certificates, as well as
24 leading-coefficient checks for $d=1,\ldots,8$.

In particular it verifies \eqref{eq:main-factor},
\eqref{eq:positive-cubic}, \eqref{eq:cycle4-discriminant},
\eqref{eq:cycle5-discriminant}, and \eqref{eq:ratio-gap}; it also checks
every polynomial listed in \cref{sec:polynomial-table}. The shifted
polynomials used for positivity have nonnegative coefficients and positive
constant term. Their positivity on a nonnegative half-line is therefore
immediate. The full coefficient arrays are written to
\texttt{data/certificates.json}.

\subsection{Independent integer-array checks}

The script \texttt{verify.py} streams rows using \eqref{eq:Srec} and
\eqref{eq:Crec}. With its default settings, it checks all $n\le200$ for
$1\le r\le5$, both subset and cycle triangles. For each of the ten
triangles there are
\[
 \sum_{n=3}^{200}(n-2)=19701
\]
positive-support interior inequalities, giving $197010$ exact strict
inequalities in total. All passed. The order-five subset lower bound
\eqref{eq:quantitative} was also checked at all $19701$ such indices.
The order-five cycle run is explicitly marked \emph{experimental} in the
output; it is not part of the proved theorem.

As a structurally independent check, the script computes
\eqref{eq:S-coefficient} and \eqref{eq:C-coefficient} by truncated
convolution of rational coefficient lists. It compares every entry with
the recurrence computation for $n\le12$, $1\le r\le8$, and both kinds:
$1456$ entries in total, all agreeing. First-column formulas, diagonal
formulas, the exact counterexamples, and the operator obstruction are
also checked.

Deb and Sokal already report finite log-concavity testing through
$n=1000$ for the cases $r=3,4,5$ \cite[Appendix~C, Tables~1--2]{DebSokal}.
Our smaller independent run is therefore not a claim to improve that
computational range. Its purpose is to audit this implementation and the
newly derived quantitative bound; the all-$n$ result is supplied by the
proof.

\subsection{Reproduction and trust boundary}

Both scripts require only Python 3.9 or later and its standard library:
\begin{lstlisting}
python certificates.py
python verify.py
\end{lstlisting}
A longer finite run is available explicitly:
\begin{lstlisting}
python verify.py --max-n 500
\end{lstlisting}
No external package installation, network request, hidden data file, or
computer algebra service is required. The assertions are implemented as
explicit checks and are not disabled by Python's \texttt{-O} flag.

The recurrence computation uses $O(N^2)$ integer-arithmetic operations and
stores $O(N)$ entries when rows are streamed, for fixed order $r$ and
maximum row $N$. These are arithmetic-operation counts, not bit-complexity
claims: the integers grow with $N$. The independent convolution check is
intentionally restricted to small rows, where exact rational arithmetic is
cheap and implementation independence is more valuable than speed.

The scripts are not a proof-assistant formalization. Their trust boundary
is ordinary Python integer/rational arithmetic and the small supplied
polynomial implementation. The paper's proofs do not logically depend on
executing the scripts: all identities, factor signs, support conditions,
and induction steps are written out.

\section{Conclusion}\label{sec:conclusion}

The chosen problem has a clean resolution: the higher-order Stirling
subset triangle has log-concave rows exactly for orders one through five.
The proof is short at its core. Normalize the recurrence, keep all terms
in the local quadratic bound, use $k\le t/(d+1)$, and factor one
discriminant. A single sixth-order counterexample plus a decreasing
factorial ratio supplies the sharp converse.

The method also settles the third- and fourth-order cycle cases and yields
a continuous interpolating family. Its failure at the fifth-order cycle
operator is genuine and explicitly certified, not merely a failure to
find a convenient factorization. Any extension to that remaining case
must exploit a stronger property of the actual preceding rows. The
present report separates that open continuation from the completed
subset-number proof.

\appendix
\section{Polynomial identities in one place}\label{sec:polynomial-table}

This appendix collects the low-degree expressions used in the proof.
The independent variable is $t$, and $H_d(t,k)$ is affine in $k$.
The notation matches the main text; the code uses \texttt{C0} and
\texttt{C1} for the constant and linear-in-$k$ parts of $H_d$.

\subsection{Pure-square coefficient}
\begin{align*}
 A_0(t)&=0,\\
 A_1(t)&=1,\\
 A_2(t)&=4(2t^2-2t-3),\\
 A_3(t)&=9(3t^4-12t^3-9t^2+42t+40),\\
 A_4(t)&=32(2t^6-18t^5+17t^4+168t^3-55t^2-834t-630).
\end{align*}

\subsection{Subset mixed coefficient and its lower bound}
\begin{align*}
 H_0(t,k)&=0,\\
 H_1(t,k)&=2,\\
 H_2(t,k)&=4(-2k+2t-1),\\
 H_3(t,k)&=6(-9kt+9k+3t^2-6t+11),\\
 H_4(t,k)&=16(-12kt^2+36kt-54k+2t^3-9t^2+43t-51).
\end{align*}
The substitution $k=t/(d+1)$ gives
\begin{align*}
 L_0(t)&=0,\qquad L_1(t)=2,\\
 L_2(t)&=\frac43(4t-3),\\
 L_3(t)&=\frac32(3t^2-15t+44),\\
 L_4(t)&=-\frac{16}{5}(2t^3+9t^2-161t+255).
\end{align*}

\subsection{Cycle mixed coefficient}
\begin{align*}
 M_0(t)&=0,\qquad M_1(t)=2,\\
 M_2(t)&=8(t-1),\\
 M_3(t)&=-18(3t-11),\\
 M_4(t)&=-32(t-2)(2t^2+4t-51).
\end{align*}

\subsection{A direct expansion check for the main discriminant}

For readers checking \eqref{eq:main-factor} without software, both sides
expand to
\begin{align*}
 4A_4-L_4^2={}&\frac1{25}\bigl(
 5376t^6-66816t^5+198528t^4+1018368t^3\\
 &\hspace{20mm}{}-7986816t^2+18351360t-18662400\bigr).
\end{align*}
This displayed coefficient form is itself included in the certificate
checks in the distributed version. The factored form is preferable for
proving its sign, because it separates the domain condition $t\ge10$
from the algebraic identity.

\section{Source and status notes}\label{sec:status}

The problem identification uses the specific version
\texttt{arXiv:2507.18959v1}. The conjecture and recurrence were checked in
both its HTML and PDF renderings. On the report date the arXiv abstract
page showed only the July 2025 version in its submission history.
Searches for the exact paper identifier, the phrase ``higher-order Stirling''
with log-concavity, and relevant arXiv author/topic combinations did not
locate a subsequent resolution. The BCC30 statement and the recent
triangular-recurrence paper \cite{Shankar} were checked separately.
This is a documented search, not proof of absence from all literature.

The mathematical claims established here can therefore be assessed
independently of any priority judgment. The result-to-source map is:
\begin{center}
\small
\renewcommand{\arraystretch}{1.25}
\begin{tabular}{>{\raggedright\arraybackslash}p{0.34\textwidth}>{\raggedright\arraybackslash}p{0.56\textwidth}}
\toprule
Source assertion & Result in this manuscript\\
\midrule
Deb--Sokal 1.4(c), row log-concavity & All stated orders $1\le r\le5$;
strictness and sharp converse also proved.\\
Deb--Sokal 1.3(c), row log-concavity & Orders $1\le r\le4$ proved;
order five not settled.\\
BCC30 Problem 30.9 & Cycle orders three and four proved;
order five not settled.\\
Other clauses of those conjectures & Not claimed as consequences of the results here.\\
\bottomrule
\end{tabular}
\end{center}

\begin{thebibliography}{9}
\bibitem{DebSokal}
Bishal Deb and Alan D. Sokal,
\emph{Higher-order Stirling cycle and subset triangles: Total positivity,
continued fractions and real-rootedness},
\href{https://arxiv.org/abs/2507.18959}{arXiv:2507.18959v1}, 2025.
See Lemma~1.1, Conjectures~1.3(c) and 1.4(c), and Appendix~C.

\bibitem{Cameron}
Peter Cameron (editor), \emph{Problems from BCC30},
\href{https://arxiv.org/abs/2409.07216}{arXiv:2409.07216v1}, 2024.
Problem~30.9, presented by Bishal Deb.

\bibitem{Sagan}
Bruce E. Sagan, \emph{Inductive and injective proofs of log concavity results},
Discrete Mathematics \textbf{68} (1988), 281--292.
\href{https://doi.org/10.1016/0012-365X(88)90120-3}{doi:10.1016/0012-365X(88)90120-3}.
See Theorem~1 and its inductive proof.

\bibitem{Shankar}
Umesh Shankar, \emph{Log-concavity of rows of triangular arrays satisfying
a certain super-recurrence},
\href{https://arxiv.org/abs/2508.12467}{arXiv:2508.12467v1}, 2025.
See Theorems~1.2--1.4 for recurrence criteria and their relationship.

\bibitem{OEIS}
The OEIS Foundation Inc., \emph{The On-Line Encyclopedia of Integer Sequences},
\href{https://oeis.org/A134991}{entry A134991: Triangle of Ward numbers},
accessed 19 September 2026.
\end{thebibliography}
\end{document}
